% Translated from: thesis/chapter_07_总结与展望.md
% Target journal: General
% Generated: 2026-04-03
\chapter{Conclusion and Future Work}
\setcounter{figure}{0}
\setcounter{equation}{0}
\renewcommand{\thefigure}{\thechapter-\arabic{figure}}
\renewcommand{\theequation}{\thechapter-\arabic{equation}}

This thesis investigates search enhancement in the context of General Game Playing (GGP), with a focus on how to construct a neural value evaluation method that can be stably integrated into Monte Carlo Tree Search (MCTS) under a unified rule-description and state-machine framework. To address the limitations of existing methods in unified representation, supervision construction, and online integration stability, this thesis proposes the Search-Distilled Residual Progress-guided Value (SDRPV) method and completes an end-to-end implementation spanning rule execution, sample construction, model training, and online evaluation.

The core idea is to organize value learning as a unified process of low-cost baseline estimation, strong-search teacher supervision, and residual correction. Following this idea, this thesis builds unified rule-execution and state-machine interfaces for \texttt{breakthrough}, \texttt{connectFour}, and \texttt{hex}; designs an SDRPV sample structure containing \texttt{s}, \texttt{b}, \texttt{q_t}, \texttt{z}, and \texttt{phi}; applies a board-token-based unified state representation and a Transformer value network for residual learning; and finally integrates the trained value model into a shared MCTS kernel, yielding a neural-value-enhanced search scheme for fixed-time online matches.

Experimental results show that SDRPV achieves stable improvements at both offline-metric and online-play levels. In offline evaluation, SDRPV outperforms baseline value \texttt{b} and the teacher-only route on \texttt{MAE}, \texttt{MSE}, and correlation metrics, indicating that the model learns effective corrections to baseline estimates. In the online baseline benchmark, SDRPV reaches an overall three-game win rate of \texttt{0.7133} against both \texttt{pure_mct} and \texttt{heuristic_mcts}, while also maintaining an overall positive edge over \texttt{baseline_token_transformer}; the ablation benchmark further shows that teacher supervision provides an effective direction for value correction and that the residual objective improves method stability. These results demonstrate that the proposed technical route can be translated into real search gains under a unified GGP framework.

The main value of this work lies in validating the feasibility of using neural value methods to serve search enhancement within a rule-driven framework, and in showing that, compared with directly fitting strong-search values, residual correction around baseline estimation is better aligned with value-modeling needs in current GGP settings. In addition, by organizing experiments with a unified state machine, a unified search skeleton, and a unified evaluation protocol, this thesis makes method behavior interpretable and comparable within clear engineering boundaries.

At the same time, this thesis has limitations. The current training data mainly come from \texttt{connectFour} self-play samples, so cross-game supervision coverage is still limited. In per-game results, SDRPV performs strongly on \texttt{breakthrough} and \texttt{connectFour}, but gains are relatively limited on \texttt{hex}. Moreover, exploration of online integration strategies between neural evaluation and MCTS remains preliminary, and there is room for further optimization of the integration mechanism.

Future work can proceed in four directions. First, introduce more games and richer stage-wise samples to improve cross-game coverage of supervision signals. Second, incorporate phase-aware modeling to enhance adaptation across different game stages. Third, continue to optimize neural-evaluation and MCTS integration, including more robust selective-calling and dynamic-budget-allocation strategies. Fourth, further verify the stability and applicability of the proposed method on larger benchmarks and more diverse GGP tasks.

Overall, this thesis addresses value correction in unified GGP settings by proposing and validating SDRPV, forming a technical route with clear supervision sources, explicit system boundaries, and reproducible experimental results. Although the current work is still at a staged level, the findings indicate that introducing search-distillation supervision and residual value correction into a rule-driven execution framework can deliver real and stable search gains under limited budgets, and provides a practical basis for follow-up research.
